\chapter{A Unifying Model}\label{chap:a_unifying_model}
We develop a unifying model of gaze that integrates informational awareness while accurately capturing the statistics of search dynamics. \cref{sec:model_setup} formulates the model in terms of FNS dynamics and introduces the search task used throughout this thesis. \cref{sec:dynamical-information_adaptation} presents the dynamical-informative kernel developed in this work, outlining its modelling motivations (\cref{subsec:motivations}) and explaining how informational kernel deposition gives rise to a regular, information-driven temporal structure in transitions between regions of interest (\cref{subsec:adaptive_kernel_dynamics}). \cref{sec:local_harmonic_oscillations} derives and quantifies the prediction from neural dynamics that localised trajectories exhibit high-frequency simple harmonic oscillations. \Cref{sec:levy_superdiffusion} examines Lévy superdiffusion, the defining characteristic of FNS dynamics, deriving the short-time, noise-driven scaling of fractional absolute moments and the monoscaling behaviour associated with the model’s self-similar superdiffusion. These dynamical predictions are compared with eye-tracking data in the next chapter.

%%% 2.1 Model setup
\section{Model setup}\label{sec:model_setup}

Having established the individual spatiotemporal signatures of FNS dynamics, we now define a standard computational framework to test these dynamics against the human behavioral data presented in \cref{chap:empirical_validations}. Because this thesis models human eye movements during visual search, we emulate the visual foraging task designed by our collaborators in \textcite{naher_primate_saccades}, which requires subjects to sequentially scan an array of dots to locate a hidden target.

\subsubsection*{Empirical protocol from Näher et al.}
In the human protocol established by \textcite{naher_primate_saccades}, trials commenced with a pre-search fixation period of 2500 to 3500~ms. Following a 100~ms gap, participants were presented with a search array of white dots, each subtending \qty{0.3}{\degree} of visual angle. Exactly 20\% of these dots functioned as true targets, and subjects were allotted a maximum search time of 5000~ms to successfully foveate a target. To investigate the influence of spatial structure, the study compared a REGULAR task---where dots were spaced identically to the centers of an invisible hexagonal tiling---with an IRREGULAR task, which randomly redistributed the same number of dots across the screen to break spatial regularity. Assuming a standard viewing distance of \qty{57.3}{\cm} (where $1^{\circ} \approx 1~\text{cm}$) and a monitor pixel density of \qty{36.22}{\px\per\cm}, spatial measurements can be directly translated between pixels and degrees using a scaling factor of \qty{36.22}{\px\per\degree}.

\subsubsection*{Computational model setup}
The empirical conditions from \textcite{naher_primate_saccades} are translated into numerical model parameters using units of spatial pixels and temporal seconds. The effective potential is constructed as a two-dimensional grid of Gaussian distributions, which act as quadratic potential wells. This simulated landscape consists of a $16 \times 16$ square grid of 256 potential wells, with 20\% of these wells randomly designated as true targets. Based on the average spatial distribution of eye movements during fixation, the standard deviation of the Gaussian wells is set to $\sigma_{\text{well}} = 3~\unit{\px}$. The average axial distance between wells is fixed at \qty{110}{\px}, which corresponds to a visual angle of approximately \qty{3.04}{\degree}. This spacing bridges the dense (\qty{2.3}{\degree}) and intermediate (\qty{4.5}{\degree}) spacings in REGULAR search conditions used in the human trials. To parameterise the spatial acuity scale of the dynamical-informative kernel $k(\cdot)$, $\sigma = \qty{80}{\px}$ was chosen to approximate human visual acuity. Visual acuity typically decays radially at a scale of roughly 2.5 degrees, which translates to 87.2 pixels assuming a viewing distance of 50 cm and a screen pixel density of 4 pixels per mm. Finally, the time constant $\tau = 1 \times 10^{-4}\,\unit{\second}$ was selected to align with the rhythmicity of human saccades. While the model yields an average inter-saccadic interval (ISI) of 360 ms, empirical data shows a saccade rate of 3.53 Hz, which corresponds to a 280 ms average ISI. Together, these elements define the 2D saliency map that acts as the explored effective potential. The full computational model whose solutions are studied throughout this work is:
\begin{align*}
    d\textbf{x}_t &= -\gamma c_{\alpha} \nabla V_{t}(\textbf{x})dt + \beta \textbf{p}_t dt + \gamma^{1/\alpha} dL_t^{\alpha},\\
    d\textbf{p}_t &= - \beta c_{\alpha} \nabla V_{t}(\textbf{x}) dt,\\
    \partial_t V_{t}(\textbf{x}) &= \frac{1}{\tau}k(\textbf{x} - \textbf{x}_t)
\end{align*}
Finally, the model parameters are calibrated to recover the distinct dynamical patterns analysed in \cref{chap:empirical_validations}. Thus, unless otherwise mentioned in this work, the baseline parameters considered for a standard experiment such as was conducted in \textcite{naher_primate_saccades}:
\begin{itemize}
    \item $\alpha = 1.2$ is selected to capture the strong search Lévy superdiffusion found in \cref{sec:search_superdiffusion} while maintaining numerical stability.
    \item $\sigma_{\text{k}} = 80~\unit{\px}$ in a kernel function $k(\textbf{x}) = \exp \{ \frac{-\|\textbf{x}\|^2}{2\sigma_{k}^2} \}$ scales the adaptive kernel to match the spatial limits of human visual acuity, as discussed above.
    \item $\tau = 2 \times 10^{-4}~\unit{\second^{-1}}$ is chosen to produce theta-rhythmic saccadic transitions, as derived in \cref{eq:theta_rhythm} and compared in \cref{sec:theta_rhythmicity_of_saccades}.
    \item $\beta = 1000~\unit{\px~ \second^{-1}}$ and $\gamma = 1000~\unit{\px^2~ \second^{-1}}$ ensures the recovery of gamma-rhythmic harmonic oscillations, placing the system well within the underdamped oscillatory regime for the \qty{3}{\px} potential wells. These are compared to data in \cref{sec:gamma_oscillations}
\end{itemize}



%%% 2.2 Dynamical-Informative adaptation
\section{Dynamical-Information adaptation}\label{sec:dynamical-information_adaptation}
The most important component of the dynamics here studied is the dynamical-informative kernel. This component is the most important extension of the work done on fractional neural sampling in works such as \textcite{chen_spatiotemporal_2022} to a full computational model of visual searching during active vision. We call the kernel dynamical-informative, as the kernel influences the model through two coupled processes: the dynamics of the FNS sampler during exploration of a landscape, which is informed by; the information gained during such exploration when considered as an active search. As such, the development of kernel dynamics is informed by previous work in statistical physics and abstract computational sampling, as well as work done informational gathering in the modelled system of visual acuity during search tasks. Importantly, a similar computational process of kernel deposition is called metadynamics in Bayesian probability landscape exploration, see \textcite{barducci_well-tempered_2008} for the introduction of metadynamics and \textcite{bussi_using_2020} for its developed use as ``well-tempered metadynamics'' in balancing sampled accuracy and ergodicity. We shall see that these kernel dynamics encode the theta-rhythmic scale of saccades explored in \cref{sec:theta_rhythmicity_of_saccades} as well as the rate of information gain explored throughout \cref{chap:functional_advantages} when considering functional roles of model dynamics in search tasks abstractly.

% We introduce a process by which the FNS walker will dynamically adapt to the explored target landscape. Dynamic adaptation is the most important addition of the original work on the FNS framework; both the circuit model in \textcite{qi_fractional_2022} and its application to visual processing in \textcite{chen_spatiotemporal_2022}.
\subsection{Motivations for the introduction of kernel adaptation}\label{subsec:motivations}
The motivations for an addition of a dynamic landscape adaptation of the model are threefold. Firstly, the quantity of information gathered by the gaze trajectory must be accounted for, as this is the fundamental purpose of a visual search. Secondly, the regularity of information-seeking saccades during visual search must be present, most importantly the recently discovered phenomena of strongly theta-rhythmic saccades explored further in \cref{sec:theta_rhythmicity_of_saccades}, as well as the well-studied phenomenon of inhibition of return (IoR) \cite{ludwig_mechanism_2009, redden_inhibition_2021}. Thirdly, any functionally informative description of search dynamics must couple these search dynamics to the current information which has been gathered complete the feedback loop of active vision \cite{Renninger2005-ee, Bruce_2009_Saliency, mcdonnell_introductory_2011, kittler_effective_2026}. 

To couple gaze dynamics to information gained continuously, we examine the structure of visual information gathering. In this work the target distribution $\pi_t(\textbf{x})$ will encode the belief in the presence a searched-for target object in 2D space, in a Bayesian sense. Thus, any information during exploration of non-target locations will cause a depletion of belief (probability mass) from the target distribution $\pi$ locally. This is equivalent in the dynamical model to raising the potential landscape $V_t(\textbf{x}) = -\ln\pi_t(\textbf{x})$ locally by some kernel centred at the gaze position whose width represents the spatial acuity of information gathering. This gives a very well-motivated definition of continuous information gathering, as classical information is exactly equated to the logarithm of probability of `surprise,' which is exactly $\pi_t(\textbf{x})$ in our search task.
 Thus, we have a well-motivated instantaneous dynamical law of kernel dynamics as the gaze is positioned at $\textbf{x}_0$:
\begin{equation}
    \partial_t V_{\text{eff}}(\textbf{x}) = \frac{1}{\tau}k(\textbf{x} - \textbf{x}_0)
\end{equation}
where $k(\cdot)$ is assumed to be some approximately radially symmetric kernel function which limits to zero for large $\textbf{x}$, and is chosen to have scaling $k(\textbf{0})=1$. $k(\cdot)$ will then have a width which correlates directly to the spatial scale of information gathering, which for the human eye is explored here in \cref{subsec:spatial_scaling}. Throughout this work, we use an isotropic Gaussian function for $k(\cdot)$, with
\begin{equation}
    k(\textbf{x}) = \exp \left( - \frac{\| \textbf{x} \|^2 }{2 \sigma_k^2} \right)
\end{equation}
\subsection{Adaptive kernel dynamics as a driver of slow rhythmic exploration}\label{subsec:adaptive_kernel_dynamics}
We can analytically predict the rate of excursions of an FNS-walker to other modes of the target distribution $\pi(\textbf{x})$. We consider excursions purely caused by the dynamically adaptive potential kernel controlled by the parameter $\tau$ which takes the analogy to the rate of information acquisition defined by some bell-shaped kernel $k(\textbf{x})$. Based purely on the value of $\tau$, and the approximate shape of $k(\textbf{x})$ we can construct a tractable transition condition as follows: The effective potential well in log-probability space becomes concave-down. Such a condition will no longer permit localisation of the FNS walker, and Lévy-driven noise will create excursions to other modes.\\

We analyse this condition by considering a single isotropic radial Gaussian target distribution centred at zero with standard deviation $\sigma_{\text{fix}}$. We simplify this problem to simply predict the average interval between excursions within a very localised potential, with excursions largely driven by deterministic kernel dynamics. It is a well known fact that Lévy noise will also introduce Poisson-distributed waiting times between frequent excursions, however analysis of this is beyond the scope of this project. This corresponds to a quadratic potential well in log-probability space. This choice is maximally general, as all local modes will be approximated to second order by a quadratic potential well, using Laplace's approximation. We also consider an adaptive kernel $k(\textbf{x})$ which is a radial Gaussian in log-probability space with a standard deviation $\sigma_{k}$. For similar reasons as the choice of potential well, such a choice yields the same results for different choices of $k(\textbf{x})$ which have the same radial curvature. Thus, using a radial coordinate $r$ which is centred at the mode $r=0$, assuming that $\sigma_{k} \gg \sigma_{\text{fix}}$ and that the walker is localised at $r \approx 0$, our effective potential function $S(r;t)$ will evolve as
\begin{align}
    \frac{\partial S(r;t)}{\partial t} &= \frac{1}{\tau} k(\textbf{x}) = \frac{1}{\tau} \exp \left( -\frac{r^2}{2 \sigma_{k}^2}\right), \\
    S(r;0) &= -\ln \pi(\textbf{x}) = \frac{r^2}{2 \sigma_{\text{fix}}^2}
\end{align}
after integrating, this yields the form effective potential after time $t$
\begin{equation}
    S(r; t) = \frac{r^2}{2 \sigma_{\text{fix}}^2} + \frac{t}{\tau} \exp \left\{ -\frac{r^2}{2 \sigma_{k}^2}\right\}
\end{equation}
Then, our excursion criterion is that at time $t_{\theta}$ at the fixation $\|x\|=0$, the radial second derivative will become negative, leaving the walker unconstrained. The second radial derivative is here
\[
 \frac{\partial^2 S(r;t)}{\partial r^2} = \frac{1}{\sigma_{\text{fix}}^2} + \frac{t}{\tau}\left[ \frac{r^2}{\sigma_k^4} - \frac{1}{\sigma_k^2} \exp \left\{ -\frac{r^2}{2 \sigma_{k}^2}\right\}\right]
\]
and we find where this second derivative is zero at the fixation point $r=0$ to find the timescale $t_{\theta}$
\[
 \frac{1}{\sigma_{\text{fix}}^2} + \frac{t_{\theta}}{\tau}\left[ 0 - \frac{1}{\sigma_k^2}  \right] = 0
\]
\begin{equation}\label{eq:theta_rhythm}
    t_{\theta} = \tau \frac{\sigma_{k}^2}{\sigma_{\text{fix}}^2}
\end{equation}
It is unsurprising that $t_{\theta}$ is proportional to $\tau$. The fact that $t_{\theta}$ is proportional to the squared ratio of the width of the adaptive kernel $\sigma_k$ and the width of the mode $\sigma_{\text{fix}}$ is of interest. This result shows that the average amount of time dwelt within a potential well will be inversely proportional to the square of the largest variance.

%%% 2.3 Local harmonic oscillations
\section{Local harmonic oscillations}\label{sec:local_harmonic_oscillations}
One predictive feature of FNS dynamics is the local oscillations in sampling within confined potentials. We can quantify such high-frequency oscillations under FNS dynamics by considering the deterministic part of our FNS equations of motion in \cref{eq:fns_system}. Considering a Gaussian probability distribution of fixation positions $\pi(\textbf{x}) \propto \exp \left( - \frac{\| \textbf{x} \|^2}{2 \sigma_{\text{fix}}^2} \right)$ and thus logarithmic gradients $\nabla \ln \pi(\textbf{x}) = -\frac{\textbf{x}}{\sigma_{\text{fix}}^2}$. Using this as our saliency landscape $\pi(\textbf{x})$ without considering perturbation to the explored potential from the slowly updating kernel dynamics, our equations of motion are
\begin{equation}
\begin{aligned}
    d\textbf{x}_t &= -\gamma c_{\alpha} \frac{\textbf{x}_t}{\sigma_{\text{fix}}^2} dt + \beta \textbf{p}_t dt + \gamma^{1/\alpha} dL_t^{\alpha}, \\
    d\textbf{p}_t &= -\beta c_{\alpha} \frac{\textbf{x}_t}{\sigma_{\text{fix}}^2} dt
\end{aligned}
\end{equation}
where we ignore the stochastic noise driving this equation, only to consider the coupled ODE which describes the deterministic part of $\textbf{x} = \textbf{x}_t$ and $ \textbf{p} = \textbf{p}_t$, where the loss of subscript is typical when moving from an SDE to an ODE
\begin{equation}
\begin{aligned}
    \frac{d \textbf{x}}{dt} &= -\gamma c_{\alpha} \frac{\textbf{x}}{\sigma_{\text{fix}}^2} + \beta \textbf{p}, \\
    \frac{d\textbf{p}}{dt} &= -\beta c_{\alpha} \frac{\textbf{x}}{\sigma_{\text{fix}}^2} 
\end{aligned}
\end{equation}
Assuming an underdamped regime where the system is periodic, corresponding to parameter range $\beta > \gamma \sqrt{c_{\alpha}} / 2 \sigma_{\text{fix}}$, the resonant ordinary frequency $f = \omega / 2 \pi$ is given by
\begin{equation}\label{eq:resonant_frequency}
    f = \frac{1}{2 \pi \sigma_{\text{fix}}} \sqrt{\beta^2 c_{\alpha} - \frac{\gamma^2 c_{\alpha}^2}{4\sigma_{\text{fix}}^2}}
\end{equation}
% The parameters $\beta$ and $\gamma$ drive the high-frequency dynamics observed during fixation intervals. $\beta = 1000$ was explicitly tuned to match the empirical microtremor fixation frequency at 65 Hz with a standard deviation of 3 pixels. Correspondingly, $\gamma = 1000$ was selected to be of a similar order of magnitude to $\beta$ to dampen these strong 65 Hz oscillations appropriately. The interaction of these parameters successfully replicates a prominent gamma-band peak centred at 65 Hz, matching the empirical eye-velocity power spectral density during stable fixations.


%%% 2.4 Lévy superdiffusion
\section{Lévy superdiffusion}\label{sec:levy_superdiffusion}
The equations of motion for our system are derived from the over-damped Langevin regime, driven by Lévy noise. Such underlying dynamics are known as Lévy flight (see \ref{subsec:levy_motion}) and are the paradigm superdiffusive system. However, both the adaptive kernel dynamics, and auxiliary momentum dynamics add complexity to the superdiffusive properties of the model. Quantifying superdiffusion in this system exploring a confined and dynamic potential landscape introduces two distinct and challenging problems: firstly, a kind of inter-modal superdiffusion is introduced through adaptive kernel dynamics which expel the walker from potential wells, and secondly a superdiffusion on small spatiotemporal scales results from near-ballistic momentum-driven harmonic oscillations.

\subsubsection*{Analysis of Lévy flight superdiffusion from the fractional Fokker-Planck equation}
Lévy processes are the canonical form of superdiffusion, yet the analysis of the precise form of statistical diffusion scaling requires close analysis. Lévy flight superdiffusion is analysed from the deterministic dynamics of the probability distribution of jump size $P(\| \Delta x \|)$ dynamics of the probability distribution. This is because the dynamics of the probability distribution is governed by the fractional Fokker Planck equation as in \cite{barkai_fractional_2001, boffi_probability_2023, qi_fractional_2022}. Considering the phase space of our coupled differential equations as $\textbf{z} = [\textbf{x}, \textbf{p}]$, the fractional Fokker-Planck equation can be analysed:
\begin{equation}\label{eq:fractional_fokker_planck}
    \frac{\partial P(\textbf{z}, t)}{\partial t} = - (\textbf{D} + \textbf{Q})\frac{\partial}{\partial \textbf{z}} \left[ b(\textbf{z}) P(\textbf{z}, t)\right] - \textbf{D} \mathcal{D}_{\textbf{z}}^{\alpha} P(\textbf{z}, t)
\end{equation}
Where $\mathcal{D}^{\alpha}_{\textbf{z}}$ represents the $\alpha$-fractional Riesz derivative, which generalises the operation of the Laplacian to fractional orders, and is defined in Fourier space as $\mathcal{D}^{\alpha}_{\textbf{z}}[f] = \mathcal{F}^{-1} \|\textbf{k}\|^{\alpha} \mathcal{F}[f]$, with the case of $\alpha = 2$ restoring the ordinary Fokker-Planck equation with $\mathcal{D}^{\alpha = 2}_{\textbf{z}} = -\nabla_{\textbf{z}}^2$. The $b(\textbf{z})$ drift term is presented in \cref{eq:full_drift}, and it's local approximation used throughout this thesis in \cref{eq:drift_approx}. The matrices in \cref{eq:fractional_fokker_planck} encode the form of the FNS dynamics given in \cref{eq:fns}, where 
$\mathbf{D} = \begin{bmatrix}
\gamma & 0\\
0 & 0
\end{bmatrix}$ 
and $\mathbf{Q} = \begin{bmatrix}
0 & -\beta\\
\beta & 0
\end{bmatrix}$ 
are the generalised matrices defining Langevin dynamical systems. $\mathbf{D}$ is the positive semi-definite matrix which defines the diffusion strength, while $\mathbf{Q}$ is the skew-symmetric matrix which defines the coupling between dynamical variables. Solutions of the dynamics in \cref{eq:fractional_fokker_planck} are nontrivial, as the fractional derivative results in a nonlinear multivariate integro-differential equation which is difficult to propagate, even numerically. This further motivates our method of analysing individual numerically integrated trajectories of \cref{eq:fns}. When considering diffusion over short timescales, the physics of Langevin systems are largely noise dominated. In our Lévy driven system, short times are dominated by $\alpha$-stable $dL^{\alpha}_t$ Lévy noise, across power-law distributed spatial scales of short jumps and nonlocal jumps to distant positions in short times. This is corroborated by the scaling of Lévy noise as $dL^{\alpha}_{\Delta t} = \Delta t^{1/\alpha} dL^{\alpha}_1$ in \cref{eq:levy_noise} is stronger than ballistic effects which scale in short times as $\Delta t$. Thus, considering the fractional Fokker-Planck equation evolving from a delta function initial condition $P(\textbf{z},0) = \delta^n(\textbf{x} -\textbf{x}_0)\delta^n(\textbf{p} - \textbf{p}_0)$, and constructing a leading-order short time solution up to $t=\Delta t$ by marginalising the auxiliary momentum $\textbf{p}$ and using the position coordinate $\Delta \textbf{x} = \textbf{x} - \textbf{x}_0$, we have:
\begin{equation}
    \partial_t P(\Delta \textbf{x},t) = -\gamma\mathcal{D}_{\textbf{x}}^{\alpha}P(\Delta \textbf{x}, \Delta t) - \nabla_{\mathbf{x}}\cdot
    \left[\left(\gamma \mathbf{b}(\Delta \textbf{x}) + \beta\,\mathbb{E}[\mathbf{p}\,|\,\Delta \textbf{x},\Delta t]\right)P(\Delta \textbf{x},\Delta t)\right]
\end{equation}
where we can truncate these terms to highest order using
\begin{equation}
    \mathbf{b}(\Delta \textbf{x}) = \mathbf{b}_0 + \mathcal{O}(\Delta \textbf{x}),\qquad 
    \mathbb{E}[\mathbf{p}|\Delta \textbf{x},\Delta t] = \mathbf{p}_0 + \mathcal{O}(\Delta t), \qquad
    \mathbf{v}_0 = \gamma\,\mathbf{b}_0 + \beta\,\mathbf{p}_0
\end{equation}
which gives
\begin{align}
	\nonumber
    \partial_t P(\Delta \textbf{x},\Delta t) &= -\gamma\mathcal{D}^{\alpha}_\textbf{x}P -\nabla_{\mathbf{x}}\cdot
    \left[\left(\textbf{v}_0 + \mathcal{O}(\Delta \textbf{x}) + \mathcal{O}(\Delta t)\right)P\right]\\
    \nonumber
    \partial_t P(\Delta \textbf{x},\Delta t) &= -\gamma\mathcal{D}^{\alpha}_\textbf{x}P - \mathbf{v}_0\cdot\nabla_{\mathbf{x}}P +  \left[\mathcal{O}(\Delta t^0) + \mathcal{O}(\Delta t^{1 - 1/\alpha})\right] \cdot P\\
    \partial_t \hat P(\mathbf{k},\Delta t) &= -\gamma\|\mathbf{k}\|^{\alpha}\hat P - i\,\mathbf{v}_0\cdot\mathbf{k}\,\hat P + \mathcal{O}(\Delta t^0)\\
    \implies \hat P(\mathbf{k},\Delta t) &= \exp\!\left[-\gamma\|\mathbf{k}\|^{\alpha}\Delta t - i\mathbf{v}_0\cdot\mathbf{k}\,\Delta t\right] [1 + \mathcal{O}(\Delta t)]
\end{align}
Showing that, to first order in $\Delta t$, the fractional Fokker-Planck under FNS dynamics in \cref{eq:fns_system} gives a stable distribution

\subsubsection*{Quantification of Lévy superdiffusion}

In calculating absolute fractional moments, we can reduce our 
\begin{align}
    \langle|\Delta x|^q\rangle &= C_q\int_0^\infty \left[ 
    1 - \mathrm{Re}\left(\hat P(k,\Delta t)\right)
    \right] k^{-1 - q} dk \\
    \langle|\Delta x|^q\rangle &= C_q\int_0^\infty \left[ 
    1 - \cos(v_0 k \, \Delta t)e^{-\gamma \Delta t \,  k^{\alpha}} \left[ 1 + \mathcal{O}(\Delta t) \right]
    \right] k^{-1 - q} dk \\
\end{align}
with 
\begin{equation}
    C_q=\frac{q(1-q)}{\cos(\frac{q\pi}{2})\Gamma(-q)}
\end{equation}
for $1 < q < \alpha$ and $q < 2$:
\begin{equation}    
\langle|\Delta x|^q\rangle =  C_q\left[\frac{\Gamma\!\left(1-\frac q\alpha\right)}{q}\,(\gamma \Delta t)^{q/\alpha}+\frac{\Gamma\!\left(\frac{2-q}{\alpha}\right)}{2\alpha}\,v_0^2\,\gamma^{-2}\;(\gamma \Delta t)^{\,2-\frac{2-q}{\alpha}} + \mathcal{O}(\Delta t ^ {1 + q/ \alpha}) \right]
\end{equation}
This analytically gives the prediction that a spectrum of fractional absolute moments will scale as $\langle \| \Delta x \| ^ q \rangle \propto \Delta t ^ {\lambda(q)}$ with $\lambda(q) = \frac{q}{\alpha}$. Thus, FNS dynamics are shown to give the same diffusion for a Lévy flight with $q < \alpha$, and $\lambda(q) = 1$ for $q > \alpha$, with an oscillatory diffusion to higher order due to $v_0$ momentum-driven terms. This piecewise superdiffusion is a result of the discontinuities in jump dynamics of a Lévy flight which lead to divergent moments, and may lead analysis to present Lévy flight behaviour as being a multifractal, despite it's monofractal scaling definition. There are a number of modelling choices which can be made to remedy this broken scaling law by taming Lévy discontinuities. Suggestions for describing a fuller monofractal superdiffusion would be to use a Lévy walk CRTW structure, as described in \cref{sec:stat_phys}, or otherwise impose an online low-pass filter to gaze trajectories, simulating more realistic muscular control of the eyes without overly altering the neural dynamical system.

%%% Summary
\subsection*{Summary}
In constructing the model under study in this thesis, this chapter presents three signature components of our unifying model of gaze dynamics: a kernel that couples regular information gathering to dynamical transitions, the high-frequency oscillations that govern fine-scale movements, and superdiffusive motion characterised by a single scaling exponent across spatial scales. In \cref{sec:model_setup} we compared the modelled setup with an empirical experiment performed in \cite{naher_primate_saccades}. In the next chapter, we assess the model’s predictive and explanatory value by comparing its predictions with the human eye-tracking data from that experiment.